\section{第 1 章 绪论}

\subsection{1.1 研究背景}

近些年来，由于数值天气预报技术、开放气象数据服务和Web应用开发技术的迅速发展，公众获取天气信息的方式也发生了明显的变化。早期的天气服务主要提供静态预报结果，即未来几天的气温、降水、风力等级等基本气象信息，而目前用户对气象产品的需要已经从原来的静态预报结果发展为现在的精细化、对比化、场景化的气象产品需求。一方面普通用户希望得到更加直观、更加连续的逐小时天气变化，从而为出行、生活提供帮助；另一方面，具有一定专业需求的用户，例如农业生产者、气象爱好者或者相关研究人员，更加关心不同模式之间存在的差异、探空层结结构、历史同类天气事件的统计特征以及预报误差的订正能力。

与此同时，开放气象数据源的丰富也给综合气象平台的创建打下了基础。以Open-Meteo为代表的开放预报接口使多模式预报结果的获取、整合变得更为方便；怀俄明大学提供的探空资料给大气稳定度分析、强对流诊断打下了基础；历史站点数据、再分析资料为气候趋势研究、极端事件识别提供长期样本。怎样把各种来源、时间跨度不同、业务逻辑各异的数据整合到一个平台上，用怎样的架构和交互方式向用户展示出来，成为了一个有实际操作价值的系统设计问题。

在此背景下，设计出一种可以满足普通用户需求，同时又具备一定专业分析能力的综合气象信息系统，有较高的应用价值和工程意义。该系统不但可以给杭州地区天气分析提供更加完整的工具链，还可以为之后扩展到更多的区域、更多的业务场景打下基础。

\subsection{1.2 研究意义}

本文的研究意义主要体现在以下几个方面。

在实际运用中，综合气象服务系统能克服单个预报展示平台所存在的缺陷。传统的天气页面只显示单个模式的结果或者很少的指标，不能反映出不同的模式之间存在的差异，也不能给探空、大气稳定度、历史同日极端性等更专业的内容提供支持。本文所实现的系统以多模式预报为主，把探空分析、历史统计、机器学习后处理整合在一个平台上，使用户可以在同一个界面里完成从短期预报浏览到深层次分析的各个步骤。

其次从工程实现角度出发，本文给出一个面向实际业务的中小型综合气象平台架构方案。系统不是简单的堆砌，而是以预报分析能力的建立为主线，把数据抓取、时间规整、站点配置统一、模型调用、分析计算和页面展示组织成比较清晰的模块结构。该种组织方式可以给同类型信息服务系统开发提供一定的借鉴。

再次从方法上来说，本文用历史同月同日的经验百分位构成EW指数来辅助识别出更具有代表性的一些极端天气事件，同时也对ECMWF/Open-Meteo预报数据进行了误差校正模型训练，将训练和运行时的站点口径、时次规则统一起来。这些工作没有追求复杂的算法革新，但是从系统实用性、结果可信度来说是有价值的。

最后在应用延伸方面，系统会把农业气象服务当作重要的扩展模块，尝试将气象分析能力同作物生长需求结合起来，创建起作物积温估算、风险提醒以及农事建议等功效，表现出综合气象平台向行业应用落地的可能。

\subsection{1.3 国内外研究与应用现状}

就国外的发展情况而言，气象服务平台已经普遍具有比较完善的在线化、可视化能力。部分业务平台可以提供多模式预报浏览、探空图显示、雷达和卫星叠加展示、专业参数分析等服务，说明气象服务正由原来的文本、表格形式的输出方式向交互式的平台转变。随着机器学习在数值预报后处理领域应用的增多，用历史预报和实况样本做偏差订正、概率预报和分类识别，已经成为提高预报质量的一个重要途径。

在国内，公共天气服务平台以及业务化气象平台近些年来也正在快速的发展，在精细化预报、短临预警、灾害监测以及行业服务等方面取得了较多的成果。但是对毕业设计或者中小型综合系统来说，常常会存在两个问题，一是平台的功能很容易停留在单个查询或者单个模块展示上，缺少多源数据融合之后的整体设计；二是即使系统包含了多个模块，也会出现功能堆砌的情况，缺少清晰的业务主线以及统一的数据组织逻辑。

因此，如何以“预报分析”为核心，兼顾历史气候、探空诊断、误差校正和农业服务等能力，构建一套结构清晰、实现完整、具备一定专业特色的综合气象信息服务系统，仍具有较强的实践研究价值。

\subsection{1.4 本文研究内容}

围绕上述目标，本文主要完成以下工作：

1. 设计并实现一个基于 Flask 的综合气象信息服务系统，构建统一的页面入口、接口组织和模块协同机制。

2. 建立多模式预报分析模块，集成 ECMWF、GFS、ICON 等模式数据，提供 72 小时精细化预报和多模式对比分析能力。

3. 建立探空分析模块，实现探空资料抓取、层结解析、稳定度参数计算和多种图形展示。

4. 建立历史气候分析模块，完成年度气候统计、趋势分析以及基于 EW 指数的代表性极端事件筛选。

5. 建立机器学习误差校正模块，对 ECMWF/Open-Meteo 预报结果进行站点统一、时次统一和模型评估。

6. 建立农业气象服务模块，将预报与历史分析能力延伸到作物生长和农事建议场景中。

需要说明的是，本课题开题阶段主要围绕综合天气展示、预报分析和历史气象等基础能力展开，农业气象服务并未作为独立模块被明确纳入正式开题边界。随着系统开发逐步推进，本文在完成核心预报分析能力的基础上，进一步扩展形成了农业气象服务链路。因而，论文正文以当前实际完成的系统为主要依据，在总体研究方向保持一致的前提下，对模块重点和结构安排作了适应性调整，即以预报分析为主线，以农业应用作为拓展内容展开论述。

\subsection{1.5 论文组织结构}

本文共分为七章。第一章为绪论，介绍研究背景、研究意义、国内外现状以及本文的主要工作。第二章介绍系统实现涉及的相关技术与理论基础。第三章对系统需求进行分析。第四章给出系统总体设计。第五章详细阐述各关键功能模块的设计与实现。第六章对系统进行测试与结果分析。第七章总结全文工作，并对后续改进方向进行展望。

